\section{Evaluation}
\label{sec:evaluation}

We evaluate four research questions---one per challenge of \S\ref{sec:intro}: \textbf{RQ1} (semantic consistency across artefacts), \textbf{RQ2} (parameter calibration and batch drift), \textbf{RQ3} (fidelity to the counterpart and divergences) and \textbf{RQ4} (service and counterpart cost). Every result carries an explicit provenance tag: \textbf{[P]} (open-source prototype, regenerable from committed artefacts), \textbf{[S1]} (frozen pre-registered Study~1, commit \texttt{9f84aad9}) or \textbf{[S2]} (frozen pre-registered Study~2, commit \texttt{0bbe06af}). Full per-scenario matrices and test breakdowns are archived in the online supplement~\citep{sag-software}.

\subsection{Subjects}

Subjects come from the framework's byte-reproducible scenario corpus~\citep{sag-software}: the 74-component \texttt{atm} scenario, a five-member ATM scale family (29--444 components), eleven generated domain scenarios (including autonomous vehicles, IoT, healthcare, telecom RAN, microservices and industrial control), two pilot scenarios excluded from the study populations (\texttt{tiny\_system}: 4 hosts, 12 apps, 2 brokers; and an integration-hub pilot), and five real-world OSS architecture topologies (\texttt{autoware\_ros2}, \texttt{edgex}, \texttt{homeassistant}, \texttt{trainticket} and \texttt{cloud\_microservices}) included in Study~2 and the perturbed-counterpart evaluation.

\subsection{RQ1: Consistency across twin artefacts}
\label{sec:consistency}

Consistency is enforced by \textbf{103 automated pytest invariants} in CI [P] covering rule-by-rule \textsc{depends\_on} parity across Rules~1--7 (15 checks), in-memory vs.\ graph-database repository parity (3), substrate-import isolation (1), oracle-role separation (8), byte-identical scenario regeneration (21), M2S/chaos/in-process parity (34) and \texttt{hybrid-verify} contracts (21; full breakdown in supplement Table~S4~\citealp{sag-software}).

These invariants caught a real defect: the simulator's re-derivation originally implemented only 1 of 7 rules, silently zeroing $I_M$ for all scored components; the fix enforced full re-derivation and distinguished \emph{not measured} from zero. Conversely, one residual divergence remains documented as visible consistency debt: when a topic is reachable both directly and via a library, the in-memory repository double-counts its Rule~1 weight while the graph-database repository (and simulator) do not (0 pairs on \texttt{atm}, 64 on \texttt{microservices} with $\Delta w \le 0.248$, 139 on \texttt{av} with $\Delta w \le 0.374$ [P]).

\textbf{Answer to RQ1.} Automated CI invariants enforce rule-by-rule semantic agreement across all artefacts, caught a major historical bug, and keep residual edge-weight debt measured and visible.

\subsection{RQ2: Calibration identifiability and perturbed-counterpart recovery}

Table~\ref{tab:identifiability} (\S\ref{sec:identifiability}) shows that on ATM, calibration recovers topic rates within 3\% median error, application service means/CVs within 1\%, and trigger policies with 100\% accuracy when discriminating outages occur, while staleness windows and burstiness under periodic workloads remain non-identifiable; broker service times, \texttt{k\_of\_n:$K$} policies and bursty MMPP workloads are recovered in unit round trips [P]. On the unperturbed Study~1 counterpart, H5 is undecidable because uncalibrated S2 already has 0.000 per-reader delivery error in seven of eight scenarios (0.018 in the eighth, unchanged after calibration; HL $0.000$, Holm n.s.) [S1].

\paragraph{Perturbed-counterpart calibration and batch drift (Docker) [P].} To test calibration and R7 batch drift when pre-deployment guesses are actually misspecified, we perturbed the Docker + Cyclone DDS counterpart on \texttt{tiny\_system} (4 containers), \texttt{atm} (8 containers) and \texttt{realworld\_\allowbreak autoware\_\allowbreak ros2} (6 containers)---shifting 40\% of topic rates by $+50\% / -35\%$ and doubling service times on the busiest host while keeping the twin's prior at declared values (40\,s runs, 5\,s warm-up, 3 repetitions, calibration from a disjoint trace; supplement Table~S6~\citealp{sag-software}):
\begin{itemize}[leftmargin=*,nosep]
  \item \emph{Batch calibration (S2 $\to$ S3):} Cuts median delivery error by 77\% on \texttt{tiny\_system} ($0.277 \to 0.064$), 96\% on \texttt{atm} ($0.053 \to 0.002$) and 87\% on \texttt{realworld\_\allowbreak autoware\_\allowbreak ros2} ($0.284 \to 0.036$), and lowers p50/p95 latency error by 10--63\% across all three topologies (e.g.\ \texttt{realworld\_\allowbreak autoware\_\allowbreak ros2} p95: $0.620 \to 0.228$).
  \item \emph{Incremental EMA (\texttt{update\_profile}, $\alpha = 0.5$):} Delivery error converges monotonically across three rolling windows on every topology (\texttt{tiny\_system}: $0.161 \to 0.091 \to 0.087$; \texttt{atm}: $0.059 \to 0.040 \to 0.029$; \texttt{realworld\_\allowbreak autoware\_\allowbreak ros2}: $0.118 \to 0.062 \to 0.038$).
  \item \emph{Drift detection and residual feedback:} \texttt{detect\_drift} achieves 1.00 TP (shifted-topic recall 0.73--1.00), but also flags unperturbed baselines ($\text{FP} = 1.00$) despite 0.28\% median baseline delivery error (0.28--2.86\% across the three topologies) because single-machine p50 latency exceeds 25\% on $>20\%$ of topics---showing that rate/delivery and latency drift require separate thresholds. Under a behavioural shift (one app changed from \emph{periodic} to \emph{all-of}), \texttt{diagnose\_fault\_residuals} correctly infers \emph{all-of} on \texttt{tiny\_system} (\texttt{A11}) and \texttt{realworld\_\allowbreak autoware\_\allowbreak ros2} (\texttt{A10}), and \emph{any-of} on \texttt{atm} (\texttt{A1}).
\end{itemize}

\textbf{Answer to RQ2.} Telemetry recovers rates, service times and observable trigger policies; under real parameter misspecification on Docker + DDS counterparts, batch calibration and rolling EMA cut delivery error by 77--96\% [P] and Study~2 confirmatorily validates delivery-error reduction across all 14 topologies (H2.3: median $0.082 \to 0.018$, HL $+0.096$, Holm $p = 0.00049$) [S2], whereas unperturbed Study~1 H5 has zero baseline error to reduce [S1]. Staleness windows and periodic-trace burstiness remain unidentifiable.

\subsection{RQ3: Fidelity of the simulators to the counterpart}
\label{sec:rq3}

\paragraph{Smoke test and a real middleware fidelity gap [P].} On \texttt{tiny\_system} in Docker (4 containers, 1\,Hz; supplement Table~S5~\citealp{sag-software}), fault-free, host-crash (\texttt{docker kill}) and broker-failover runs agree exactly on stalled readers ($J = 1.0$), exact DDS \texttt{liveliness\_\allowbreak alive\_\allowbreak count} across all 15 readers, and per-reader delivery within 5.9\% (p95 latency error 11--18\%, within 35\% tolerance; p50 14\% fault-free, 34--40\% under faults). Next, a 10\,s switch-crash partition on a two-host chain (with network gating and CPU-bound service matching 21.3\,ms vs.\ 22.5\,ms exponential median) stopped deliveries within 0.2\,s as simulated, but resumed \textbf{$\approx 6$\,s after the partition healed} whereas $\Idyn$ reconnects instantaneously: during partitions longer than the 1\,s participant lease, Cyclone DDS expires remote participants and must rediscover them. This independent-row finding (Table~\ref{tab:independence}) is invisible to any self-consistent simulator.

\paragraph{Pre-registered fidelity studies [S1] and [S2].} In Study~1 [S1], all 2,738 counterpart runs across $n = 8$ scenarios at $\kappa = 1$ passed validity on the first attempt (0.000 pre-fault rate error; peak CPU p95 38--62\%), simulator arms ran 10,950 times, and the instrument check passed ($\rho(\mathrm{N}) = 0.001$). In Study~2 [S2], all 14 topologies (13 at $\kappa = 1$, plus \texttt{cloud\_microservices} at $\kappa = 0.74$) completed cleanly ($\approx 18.0$\,h, $\approx 853$ Docker runs). Table~\ref{tab:study} summarises both studies (per-scenario matrices in supplement Tables~S7--S8~\citealp{sag-software}).

\begin{table}[t]
\centering
\caption{Pre-registered fidelity study results: Study~1 [S1] ($n = 8$, $\kappa = 1$, $\delta = 0.05$) and Study~2 [S2] ($n = 13$--$14$, independent rows). HL: Hodges--Lehmann median [95\% CI]. Holm-corrected secondary tests are two-sided.}
\label{tab:study}
\small
\begin{tabularx}{\textwidth}{@{}>{\raggedright\arraybackslash}p{0.26\textwidth} >{\raggedright\arraybackslash}Y >{\raggedright\arraybackslash}p{0.27\textwidth} >{\raggedright\arraybackslash}p{0.14\textwidth}@{}}
\toprule
Quantity & Statistic & Value & Decision \\
\midrule
\multicolumn{4}{@{}l}{\emph{Study~1 [S1] (commit \texttt{9f84aad9}, $n = 8$, O-open publishers)}} \\
H1: S2 vs.\ S0 (publishers) & HL [CI]; rank-biserial; $p_{\text{sup}}$; $p_{\text{TOST}}$ & $+0.014$ [$-0.002$, $0.094$]; 0.56; 0.098; 0.098$^{a}$ & \textbf{inconclusive} \\
H2--H5 (secondary, Holm) & H2 S0 vs.\ 0.70 ($p = 0.016$); H3 S2 vs.\ R ($0.84$); H4 S0 vs.\ R ($0.078$); H5 cal ($0.50$) & all arms \& R at ceiling ($\rho \approx 0.91$--$0.93$ vs.\ $C_s = 0.93$) & all n.s.\ (Holm) \\
Instrument check (arm N) & median $\rho_s(\mathrm{N})$ & 0.001 & within $\pm 0.10$ \\
\midrule
\multicolumn{4}{@{}l}{\emph{Study~2 [S2] (commit \texttt{0bbe06af}, $n = 13$--$14$, independent rows)}} \\
H2.1: p50 latency S0 $-$ S2 & $n = 13$ ($\kappa = 1$); HL [CI]; 1-sided $p$ & $+0.135$ [$+0.071$, $+0.225$]; $p = 0.00061$ & \textbf{better} \\
H2.2a: Infra $\rho - 0.85$ & $n = 14$; median $\rho$ (range); Holm $p$ & $0.979$ ($0.916$--$1.000$); $p_{\text{H}} = 0.00049$ & \textbf{significant} \\
H2.2b: Infra stall $J - 0.85$ & $n = 14$; median $J$; Holm $p$ & $1.000$ (all 14/14); $p_{\text{H}} = 0.00049$ & \textbf{significant} \\
H2.3: Perturbed cal S2 $-$ S3 & $n = 14$; HL [CI]; Holm $p$ & $+0.096$ [$+0.034$, $+0.182$]; $p_{\text{H}} = 0.00049$ & \textbf{significant} \\
H2.4: Omission vs.\ crash ref & $n = 14$; HL [CI]; raw $p$ & $+0.0003$ [$-0.0006$, $+0.0021$]; $p = 0.542$ & n.s.\ (Holm) \\
\bottomrule
\multicolumn{4}{@{}p{\textwidth}@{}}{\footnotesize $^{a}$Both $p$-values equal 25/256; a tie-aware permutation test gives $p_{\text{sup}} = 0.078$ (same decision).}\\
\end{tabularx}
\end{table}

\paragraph{Interpreting Study~1 and Study~2 together.} In Study~1, six scenarios differ by $|d_s| \le 0.02$, \texttt{logistics\_\allowbreak fleet} favours S0 by 0.020, and noisy \texttt{hub\_\allowbreak and\_\allowbreak spoke} ($C_s = 0.66$) favours S2 by 0.177, leaving H1 inconclusive. Every arm (S0--S3 and closed-form R) sits at the noise ceiling ($\rho \approx 0.91$--$0.93$ vs.\ $C_s = 0.93$): under O-open, a periodic publisher crash removes only direct-reader deliveries, which R computes directly from the shared plan (dependent rows of Table~\ref{tab:independence}). H2 clears 0.70 in 7/8 scenarios, but at $n = 8$ one loss floors two-sided $p$ at $0.016 > 0.0125$.

Study~2 confirms what Study~1's descriptive analyses suggested once the contrast moves to independent rows across 14 topologies (including five real-world OSS systems): (1)~host stations (S2) significantly lower steady-state p50 latency error from median 0.376 (S0) to 0.258 (11/13 wins; p95 $\le 0.30$ on all 13 full-rate scenarios, well inside the 35\% tolerance); (2)~on broker and host crashes (247 targets total), S2 ranks impact at median $\rho = 0.979$ (\texttt{autoware\_ros2} $= 1.000$) and achieves exact stall-set agreement ($J = 1.000$) on all 14 scenarios (after fixing Study~1's \texttt{healthcare} \texttt{A4} prime-second 59\,s silence-window artefact on \texttt{main}, while Study~1's frozen 0.28 stands); (3)~under parameter misspecification, S3 reduces delivery error on all 14/14 scenarios (median $0.082 \to 0.018$, and $0.161$--$0.334 \to 0.003$--$0.047$ on the five OSS topologies); and (4)~for top-3 publisher omissions ($p = 0.5$ over 10\,s, H2.4), small impact magnitudes leave both S2 and the crash reference with tiny absolute errors ($\approx 0.001$--$0.028$, with the clearest separation on \texttt{atm}: $0.0005$ vs.\ $0.020$) and no significant difference ($p = 0.542$). Study~1's descriptive DDS-status checks (40 faults: \texttt{liveliness\_alive\_count} and \texttt{offered\_deadline\_missed} $J = 0.98$; \texttt{requested\_deadline\_missed} 0.71; \texttt{cache\_\allowbreak replaced\_\allowbreak dropped\_\allowbreak samples} 0.50), sink controls ($I_{\mathrm{obs}} \in [0.002, 0.038]$) and O-react sensitivity ($\le 0.06$) are archived in the supplement~\citep{sag-software}.

\textbf{Answer to RQ3.} The counterpart confirms smoke-test stall sets and liveliness and exposes a 6\,s DDS post-partition rediscovery gap [P]. Study~1 shows open-loop publisher ranking is inconclusive and tied with closed-form R at the noise ceiling [S1]. Study~2 confirmatorily establishes that added model detail significantly improves steady-state latency (H2.1), reproduces broker/host crash rankings and exact stall sets (H2.2a/b), and makes S2M calibration effective under parameter misspecification (H2.3), while gray-failure omission error vs.\ a crash reference is n.s.\ (H2.4) [S2]. All counterpart evidence concerns $\Idyn$ alone.

\subsection{RQ4: Cost of the twin's services and counterpart}

Table~\ref{tab:cost} shows a five-order-of-magnitude cost ladder matching continuous-delivery stages: sub-second graph gating and graph simulators fit every commit; seconds-per-fault $\Idyn$ fits candidate screening; hybrid two-tiered verification scales with shortlist size rather than $|V|$---finishing in 54.2\,s in-process on \texttt{tiny\_system} and \textbf{107.4\,min measured} in Docker + DDS on \texttt{atm} for two fault probes ($\approx 48$--54\,min/probe, a 4--6$\times$ speedup over exhaustive sweeps, achieving 1.000 baseline health-state agreement across all 66 entities, Spearman $\rho = 1.000$ across the 5 primary \texttt{crash} targets [top-$K$ Jaccard 1.0], and exact rank agreement across the 3 non-host \texttt{omission} targets [\texttt{i\_flow} MAE 0.0008, stall-set $J = 1.0$; primary-only $\rho$ omitted at $n = 2$] while transparently skipping the 5 unsupported host-level omissions in \texttt{caveats}); and exhaustive counterpart sweeps ($\approx 18.0$\,h for Study~2; $\approx 56$\,h for Study~1; 10,950 simulator runs took 69\,min on 16 cores) fit periodic twin validation.

\begin{table}[t]
\centering
\caption{Cost ladder of twin services and counterpart: [P] except the last row [S1]/[S2]. Treat timings as orders of magnitude.}
\label{tab:cost}
\small
\begin{tabularx}{\textwidth}{@{}l Y l@{}}
\toprule
Artefact & Workload & Cost \\
\midrule
Change gate (analysis + checks) & ATM scale family, 29 / 74 / 148 / 296 / 444 components & 0.01--1.01\,s \\
Graph simulators ($I^*, \Icomp$) & 5-seed / exhaustive sweep, \texttt{atm} (74) / \texttt{av} (152) & 0.06--0.45\,s \\
Queue-flow simulator ($\Idyn$) & per fault (\texttt{atm}/\texttt{av}); 5-seed whole corpus & 1.5--4.3\,s (12.7\,CPU-h) \\
Hybrid verification (in-process) & \texttt{tiny\_system}: baseline + 17 Tier-1 $\times$ 2 seeds + 2 probes & 54.2\,s \\
Hybrid verification (Docker + DDS) & \texttt{atm}: baseline + 39 Tier-1 $\times$ 3 seeds + 8 targets $\times$ 3 reps $\times$ 2 probes & 107.4\,min ($\approx$48--54\,min/probe) \\
Counterpart (Docker + DDS) & single run; Study~2 (14 scenarios, $\approx 853$ runs) [S2]; Study~1 (8 scenarios, 2,738 runs) [S1] & $\approx$71\,s; $\approx$18.0\,h; $\approx$56\,h \\
\bottomrule
\end{tabularx}
\end{table}

\textbf{Answer to RQ4.} Matching each service to its rung on the five-order-of-magnitude cost ladder makes continuous architectural verification practical without running exhaustive container sweeps per change.
